\question[5]
\enquote{Datensicherheit schützt Daten vor unbefugtem Zugriff -- ist das dann nicht dasselbe wie Datenschutz?} \\
a. Erläutere den Kernunterschied zwischen \textit{Datensicherheit} und \textit{Datenschutz}. Gehe darauf ein,
\textit{was} jeweils geschützt wird und \textit{welche Daten} jeweils betroffen sind. \\
b. Begründe die Aussage: \enquote{Datensicherheit ist eine Voraussetzung für Datenschutz, aber allein nicht
ausreichend.} \\
c. Ordne jeder Situation zu, ob \textit{Datensicherheit}, \textit{Datenschutz} oder \textit{beides} betroffen ist: \\
\hspace*{1em} (1) Eine Firma sichert die Baupläne eines neuen Produkts gegen Diebstahl. \\
\hspace*{1em} (2) Ein Online-Shop schützt die gespeicherten Namen und Adressen seiner Kunden vor Hackern
und verwendet sie nur zur Bestellabwicklung. \\
\hspace*{1em} (3) Eine Schule speichert die Fingerabdrücke aller Schülerinnen und Schüler in einer
hervorragend abgesicherten Datenbank.

\begin{solutionbox}{7cm}
a. \textit{Datensicherheit} schützt \textbf{Daten} (alle Daten) vor Verlust, Manipulation und unbefugtem
Zugriff -- im Mittelpunkt stehen technische und organisatorische Maßnahmen. \textit{Datenschutz} schützt
\textbf{Menschen und ihre Persönlichkeitsrechte} und betrifft nur \textbf{personenbezogene} Daten -- im
Mittelpunkt steht die rechtmäßige Erhebung, Verarbeitung und Nutzung. Datensicherheit fragt
\enquote{Sind die Daten geschützt?}, Datenschutz fragt \enquote{Dürfen diese Daten überhaupt erhoben
und verwendet werden?} \\
b. Sind personenbezogene Daten nicht ausreichend gesichert, können sie abfließen oder manipuliert werden
-- Datenschutz ist dann nicht mehr gewährleistet. Umgekehrt kann eine perfekt gesicherte Datenbank
trotzdem gegen den Datenschutz verstoßen, wenn die Daten ohne Rechtsgrundlage bzw. Einwilligung
erhoben oder weitergegeben werden. \\
c. (1) nur Datensicherheit (keine personenbezogenen Daten). \quad
(2) beides (Schutz vor Hackern = Datensicherheit; Zweckbindung = Datenschutz). \quad
(3) Datensicherheit erfüllt, Datenschutz möglicherweise verletzt, da die Erhebung der Fingerabdrücke
evtl. gar nicht zulässig ist.
\end{solutionbox}
